% GENERATED FILE. Edit canonical lessons, then run npm run generate:books.
% canonical-source-hash: fnv1a64:51a83aba61a2bb64
% canonical-lessons: LA-C21-volup-est-convenisse, LA-C21-meeting-phrase-context

\chapter{Nice to Meet You}
\label{ch:la-nice-to-meet-you}

This chapter is generated from the canonical micro-lessons used by Language Ladder.
Each section stays independently resumable and preserves the authored prerequisite order.

\section[volup est convēnisse]{volup est convēnisse — the phrase Latin uses instead of a fixed idiom}

\label{lesson:LA-C21-volup-est-convenisse}



\begin{quote}
\textbf{Warm-up.} {[}PAUSE 2s{]} Every modern language in this course has one fixed phrase for "nice to meet you." Latin, honestly, does \textbf{not} — but it has something better: a real line from a Roman comedy that does the job.
\end{quote}

\begin{culture}
\textbf{Plautus}'s comedy \emph{Miles Gloriosus} (Act 2, Scene 3) gives us a real line that does exactly this work: \textbf{"Tē, {[}name{]}, volup est convēnisse"} — literally "\textbf{you, {[}name{]}, {[}it{]} is-a-pleasure to-have-met}," or naturally, "\textbf{it's a pleasure to have met you, {[}name{]}.}" Break it down:

\begin{itemize}
\raggedright
  \item \textbf{volup} — an old, indeclinable word for "\textbf{pleasant, pleasing}" (functioning here the way \emph{iuvat}, "it pleases," would)
  \item \textbf{est} — "{[}it{]} \textbf{is}"
  \item \textbf{convēnisse} — the \textbf{perfect infinitive} of \emph{convenīre}, "to come together, to meet" — "\textbf{to have met}"
\end{itemize}

Put together: "\textbf{to have met {[}you{]} is a pleasure.}" This is a genuine, attested line — not a modern reconstruction.
\end{culture}

\subsection*{Guided Practice}
{[}PAUSE 1s{]}

\begin{itemize}
\raggedright
  \item {[}YOU SAY: "volup est convēnisse" — "it's a pleasure to have met you"{]}
  \item {[}YOU SAY: the pieces — volup (pleasant) + est (is) + convēnisse (to have met){]}
  \item {[}YOU SAY: the grammar — convēnisse is a perfect infinitive{]}
\end{itemize}

\begin{tcolorbox}[breakable,colback=teal!4,colframe=teal!35!black,title={Wrap-up recall}]
{[}PAUSE 3s{]} Does Classical Latin have one single fixed idiom for "nice to meet you"? (\textbf{No} — Romans expressed the sentiment according to context.) What real, attested phrase can do the job, and where does it come from? (\textbf{"Volup est convēnisse"}, literally "to have met {[}you{]} is a pleasure" — \textbf{Plautus}, \emph{Miles Gloriosus}.) What kind of infinitive is \textbf{convēnisse}? (\textbf{Perfect} — "to have met.")
\end{tcolorbox}
\section[tē volup est convēnisse]{Useful phrase, honest context}

\label{lesson:LA-C21-meeting-phrase-context}



\begin{quote}
\textbf{Warm-up.} {[}PAUSE 2s{]} \textbf{Volup est convēnisse} is genuinely Roman. Now place it in its scene so you know exactly what that evidence does—and does not—prove.
\end{quote}

\begin{culture}
In \emph{Miles Gloriosus}, Sceledrus and Palaestrio are fellow slaves who already know one another. The line welcomes a \textbf{chance encounter}. It does not document a fixed formula that Romans always used when introduced to a stranger for the first time.

The phrase still expresses the right sentiment: "to have met you is a pleasure." Use it as an attested way to do that conversational work, while remembering that Classical Latin has no single fixed counterpart to Spanish \emph{encantado} or German \emph{freut mich}.
\end{culture}

\begin{culture}
Modern spoken-Latin guides may teach \textbf{mihi pergrātum est tē convenīre}: "it is very pleasing to me to meet you." Its grammar is sound, but the formula does not occur in surviving classical texts. It is useful modern Latin, not direct evidence of Roman conversational convention.
\end{culture}

\subsection*{Guided Practice}
{[}PAUSE 1s{]}

\begin{itemize}
\raggedright
  \item {[}YOU SAY: "volup est convēnisse" — classically attested{]}
  \item {[}YOU SAY: its scene — a chance encounter between people already acquainted{]}
  \item {[}YOU SAY: "mihi pergrātum est tē convenīre" — grammatical, but modern{]}
\end{itemize}

\begin{tcolorbox}[breakable,colback=teal!4,colframe=teal!35!black,title={Wrap-up recall}]
{[}PAUSE 3s{]} Is the Plautine line a first-time introduction? (\textbf{No — it is a chance encounter between acquaintances.}) Is it still genuinely attested? (\textbf{Yes.}) Is \textbf{mihi pergrātum est tē convenīre} bad grammar? (\textbf{No.}) Is it classically attested? (\textbf{No — it is a modern spoken-Latin formula.})
\end{tcolorbox}
